\subsection{Dataset}
The first question in the evaluation is whether the implemented SSS SLAM data-processing pipeline can work on real underwater sensor data, not only on idealized or isolated test cases. The dataset is therefore important because it defines the conditions under which the state estimator, map generation, feature extraction, and performance tests are evaluated. A useful dataset must contain realistic vehicle motion, real sonar noise, real sensor timing, and enough environmental structure to test whether stable maps and repeatable landmarks can be produced.
\\ \\
The experiments in this chapter use datasets recorded during previous missions of NTNU's Light Autonomous Underwater Vehicle (LAUV) Fridtjof in the Trondheim fjord \cite{LAUV_Fridtjof}. These datasets contain synchronized measurements from the onboard navigation sensors and the side scan sonar system, providing a realistic test environment for evaluating the SLAM data-processing pipeline.
\\ \\
It is important to note that the datasets were originally recorded for general navigation and testing missions rather than specifically for this thesis. During these missions the LAUV operated with its own onboard navigation stack, including state estimation, control, and mission management components configured for normal vehicle operation. These onboard algorithms produce processed navigation outputs that are included in the recorded logs. However, these processed signals were not used directly by the SSS SLAM pipeline developed in this work, as they are tuned for operational robustness rather than high precision SLAM evaluation.
\\ \\
Instead, the SSS SLAM system presented in this thesis operates on the raw sensor measurements recorded during the missions. In particular, the raw side scan sonar data, inertial measurements, and other primary sensor streams were used as the main inputs to the processing pipeline. This makes it possible to evaluate the developed methods independently, including whether the state estimate is suitable for map generation, whether the generated sonar maps remain geometrically stable, and whether the extracted landmarks are useful for later SLAM processing.
\\ \\
The processed navigation outputs generated by the LAUV during the original missions were used only as a reference for sanity checking and benchmarking purposes. They provide a useful baseline for verifying that the estimated trajectories and system behaviour remain physically reasonable, but they do not influence the SLAM data-processing pipeline itself. This ensures that the results presented in this thesis reflect the performance of the implemented pipeline rather than the vehicles original onboard navigation system.
